% !TEX root = ../../main.tex

\section{非功能性需求分析}

\subsection{性能需求}

系统应能够满足普通用户在移动端训练和识别过程中的基本响应要求。
在普通业务操作中，如登录、资源列表加载、个人信息查看和后台管理查询等，系统应尽量保持较短的响应时间，避免用户长时间等待。
在实时手势训练场景中，系统应尽可能保证识别结果反馈的连续性，使用户能够在练习过程中及时看到识别结果。
在句子视频识别场景中，由于视频上传和识别分析耗时相对较长，系统应提供明确的等待状态提示，避免用户误认为操作失败。

对于后台管理端，系统应支持管理员在资源管理、样本管理和模型版本管理页面进行分页查询、条件筛选和状态更新。
对于模型训练和句子视频识别这类耗时任务，系统应避免阻塞普通页面操作，并通过状态提示或结果页面向用户反馈处理进度和处理结果。

\subsection{可靠性需求}

系统应具备基本的异常处理和容错能力。
在移动端网络异常、识别服务暂时不可用、视频识别失败或登录状态过期等情况下，系统应向用户给出明确提示，并允许用户重新操作。
对于后台管理端，资源新增、修改、删除、模型发布等操作应保证数据一致性，避免出现资源记录存在但文件缺失、模型版本状态错误或发布状态不明确等问题。

在实时识别和句子视频识别场景中，系统应能够处理识别失败、低置信度结果、未检测到有效动作等情况。
识别结果不稳定时，系统应以提示或结果说明的形式告知用户，而不是直接给出误导性结论。
对于样本和模型版本等关键数据，系统应保留必要的状态字段和记录信息，便于后续排查和维护。

\subsection{易用性需求}

系统移动端面向普通用户，应保证界面入口清晰、操作步骤简单、反馈信息明确。
用户在使用系统时，应能够较容易地区分首页、训练和个人中心等模块，并能够通过较少步骤进入手势训练或句子视频识别功能。
训练页面应突出动作示范、开始练习和识别结果展示，避免过多技术细节影响用户理解。

后台管理端面向管理员，应采用表格、表单、筛选条件和状态标签等常见管理界面形式，降低资源维护和模型管理的操作成本。
管理员在进行手势分类、手势资源、样本和模型版本管理时，应能够快速理解各字段含义和操作结果。
对于模型训练、模型发布等重要操作，系统应提供必要的确认提示和结果反馈。

\subsection{安全性需求}

系统应对普通用户和管理员分别进行身份认证。
普通用户登录后才能访问移动端主要功能，管理员登录后才能访问后台管理端。
系统应避免未登录用户直接进入主界面或后台页面，防止未授权访问。
登录态应通过 token 等方式进行管理，并在退出登录或登录过期后失效。

对于文件上传、头像上传、视频识别和模型文件管理等功能，系统应限制访问权限，避免普通用户越权访问管理端资源或管理员误操作用户侧数据。
后台管理端涉及样本管理、模型训练和模型发布等关键功能，应仅允许管理员角色操作，以保证识别模型和训练资源的安全性。

\subsection{可维护性与扩展性需求}

HearBridge 系统由移动端、服务端、后台管理端和 Python 识别服务组成，各部分应保持相对清晰的职责边界。
移动端主要负责用户交互和结果展示，服务端主要负责业务编排和数据管理，后台管理端主要负责资源维护和模型管理，Python 识别服务主要负责视觉识别和模型训练。
这样的分工有助于后续单独维护某一端功能，减少模块之间的耦合。

从扩展性角度看，系统后续可能继续扩充手语资源数量、增加样本规模、优化识别模型、扩展句子视频识别词表，并增加更丰富的训练反馈方式。
因此，系统在资源管理、样本管理和模型版本管理方面应预留持续扩展能力。
对于识别服务，系统应支持模型版本切换和模型重载，以便在不大幅修改移动端业务逻辑的情况下更新识别能力。
